\documentclass{beamer}
\usepackage[utf8]{inputenc}
\usepackage[spanish]{babel}
\usepackage{amsmath, amssymb}

\usetheme{Madrid}
\usecolortheme{dolphin}
\setbeamertemplate{navigation symbols}{}
\setbeamertemplate{footline}[frame number]

\title[Puente Browniano]{\textbf{Puente Browniano}}
\subtitle{Procesos Estocásticos}
\author{Dylan Gabriel Albor Saucedo}
\date{\today}

\begin{document}

\begin{frame}
    \titlepage
\end{frame}

%------------------------------------------------
\begin{frame}{Contexto histórico}
En la década de 1930, Kolmogorov y Smirnov desarrollaron pruebas estadísticas destinadas a determinar si una muestra de datos proviene de una distribución teórica específica.

\vspace{0.3cm}
Sin embargo, las demostraciones originales de estos resultados eran considerablemente complejas desde el punto de vista analítico.

\vspace{0.3cm}
En 1949, Doob estableció una conexión fundamental al proponer que las distribuciones empíricas convergen hacia un puente browniano.

\vspace{0.3cm}
Esta relación permitió simplificar de manera significativa el estudio de pruebas de bondad de ajuste.
\end{frame}

%------------------------------------------------
\begin{frame}{Definición del puente browniano}
\[
\{W(t) \mid W(T)=a\}
\]

El puente browniano se define como un movimiento browniano condicionado a alcanzar un valor fijo $a$ en un tiempo final $T$.

\vspace{0.3cm}
Este tipo de proceso conserva la naturaleza aleatoria del movimiento browniano, pero incorpora una restricción global en su comportamiento.
\end{frame}

%------------------------------------------------
\begin{frame}{Construcción del proceso}
\[
B_t = W_t - \frac{t}{T}(W_T - a)
\]

Esta expresión proporciona una construcción explícita del puente browniano a partir de un movimiento browniano estándar.

\vspace{0.3cm}
El término $(W_T - a)$ representa la desviación final, mientras que $\frac{t}{T}$ distribuye dicha corrección a lo largo del intervalo temporal.

\vspace{0.3cm}
De esta manera, el proceso es ajustado progresivamente para satisfacer la condición final.
\end{frame}

%------------------------------------------------
\begin{frame}{Verificación de la condición final}
Se verifica que el proceso cumple la condición impuesta evaluando en $t=T$:

\[
B_T = W_T - \frac{T}{T}(W_T - a)
\]

\[
= W_T - (W_T - a)
\]

\[
= a
\]

Esto confirma que la construcción garantiza el valor final deseado.
\end{frame}

%------------------------------------------------
\begin{frame}{Función de media}
\[
\mathbb{E}[B_t]
\]

Partimos de la definición:
\[
B_t = W_t - \frac{t}{T}(W_T - a)
\]

Aplicando esperanza:
\[
\mathbb{E}[B_t] = \mathbb{E}[W_t] - \frac{t}{T}(\mathbb{E}[W_T] - a)
\]

Como $\mathbb{E}[W_t]=0$ y $\mathbb{E}[W_T]=0$:
\[
\mathbb{E}[B_t] = \frac{t}{T}a
\]

Este resultado indica que el proceso tiene comportamiento promedio lineal.

\vspace{0.3cm}
Además, esta expresión permite entender cómo evoluciona en promedio el proceso y sirve como base para analizar su comportamiento global.
\end{frame}

%------------------------------------------------
\begin{frame}{Función de covarianza}
\[
\mathrm{Cov}(B_s, B_t)
\]

La covarianza describe la dependencia entre los valores del proceso en distintos instantes de tiempo.

\vspace{0.3cm}
A diferencia del movimiento browniano estándar, el puente browniano presenta dependencia global debido a la condición final.
\end{frame}

%------------------------------------------------
\begin{frame}{Cálculo de la covarianza}
En el caso estándar ($T=1$, $a=0$):
\[
B_t = W_t - tW_1
\]

Queremos calcular:
\[
\mathrm{Cov}(B_s, B_t)
\]

Sustituyendo:
\[
\mathrm{Cov}(W_s - sW_1, W_t - tW_1)
\]

Expandiendo:
\[
= \mathrm{Cov}(W_s, W_t) - t\mathrm{Cov}(W_s, W_1)
\]
\[
- s\mathrm{Cov}(W_1, W_t) + st\mathrm{Var}(W_1)
\]

Usando propiedades del browniano:
\[
= s - ts - st + st
\]

Finalmente:
\[
\mathrm{Cov}(B_s, B_t) = s(1 - t)
\]

Este resultado describe la dependencia temporal del proceso.
\end{frame}

%------------------------------------------------
\begin{frame}{Representación dinámica}
\[
dB_t = -\frac{B_t}{T - t}dt + dW_t
\]

El proceso también puede describirse mediante una ecuación diferencial estocástica.

\vspace{0.3cm}
El término $dW_t$ introduce aleatoriedad, mientras que el término determinista actúa como una fuerza de corrección que garantiza la condición final.
\end{frame}

%------------------------------------------------
\begin{frame}{Aplicación estadística}
El puente browniano desempeña un papel fundamental en la prueba de Kolmogorov-Smirnov.

\vspace{0.3cm}
En el límite, el error entre la distribución empírica y la teórica converge hacia este proceso.

\vspace{0.3cm}
Esto permite analizar el comportamiento del estadístico de prueba de manera más tractable.
\end{frame}

%------------------------------------------------
\begin{frame}{Conclusión}
\begin{itemize}
    \item El puente browniano es un proceso condicionado con propiedades bien definidas.
    \item Su estructura probabilística se describe mediante su media y covarianza.
    \item Posee una interpretación dinámica y aplicaciones relevantes en estadística.
\end{itemize}
\end{frame}

\end{document}
